\chapter{Discussion}
\label{ch:discussion}

I believe Saturn's most immediate implication concerns institutional AI provisioning.  Bassignana et al.~\cite{bassignana2025aigap} documented the consequence of the not providing AI services to students --- students from lower socioeconomic backgrounds use AI tools less frequently and benefit less from them. Gabriel~\cite{gabriel2024educational} observed the same disparity in educational outcomes. Saturn introduces a third option: an institution deploys a Saturn beacon and funds a token budget with a cloud provider; every device that joins the network discovers the service automatically, and students pay nothing upfront. The institution pays only for tokens consumed, not seats licensed.

Shared token provisioning is not without risk. A budget split across an entire campus invites resource exhaustion: a small number of heavy users could drain funds intended for thousands. Institutions would need usage caps, per-student quotas, or tiered access policies. These are governance mechanisms Saturn does not provide and fall outside the protocol's scope. I raise this not to diminish the scenario but to acknowledge that provisioning AI as a shared resource introduces allocation problems that shared printers never faced, because inference costs scale with usage in ways that printer toner does not.

Saturn's contribution to AI democratization is structural, not incremental. Costa et al.~\cite{costa2024democratization} identified accessibility and usability as the two core dimensions of AI democratization --- the ability to reach AI tools without financial or technical barriers, and the effort required to operate them once reached. The cognitive walkthrough (\argref{A13}) found that Saturn reduces end-user configuration from seven steps to zero, suggesting progress on both dimensions: no financial barrier because the institution funds the token budget, and no usability friction because discovery is automatic. Whether Saturn fully satisfies Costa et al.'s criteria requires empirical validation with real users, but eliminating all end-user steps is a stronger signal than reducing them would be.

The passive discovery model is what I find most compelling about Saturn's potential. The current subscription model requires users to commit money before evaluating whether AI is useful to them. A user on a Saturn-enabled network encounters AI capabilities through applications that consume them automatically. The user may not even recognize that AI is involved --- they simply notice that their media player now translates subtitles, or their note-taking app now summarizes documents. That invisibility lowers the threshold for first contact with AI services and could bring inference to people who would never have sought it out. But invisibility cuts both ways. If users consume AI-generated output without knowing it, they lose the opportunity to evaluate its reliability, question its biases, or opt out. Applications that integrate Saturn bear a transparency obligation the protocol itself does not enforce.

\section{Security Trade-offs of Broadcast Discovery}
\label{sec:security-tradeoffs}

Saturn deliberately trades verification for accessibility. Requiring per-device authentication would reintroduce the credential burden Saturn exists to eliminate, so I chose to document and bound the exposure instead. The question is whether that trade-off is defensible --- whether the security surface of broadcast discovery is understood well enough to manage.

\subsection{What Saturn Exposes}

Saturn's TXT records broadcast eleven metadata fields to every device on the local network segment. In the base configuration --- backend and clients on the same network --- none contain credential material. Fields like \texttt{api\_base} and \texttt{api\_type} reveal the endpoint URL and backend protocol, but neither is a secret. Cloud deployments add exactly one credential field: \texttt{ephemeral\_key}, a time-limited API key with a ten-minute maximum lifetime and five-minute rotation interval (Section~\ref{sec:ephemeral-keys}). This is the only secret Saturn broadcasts, and it is short-lived by design. All fields are visible to any device on the local segment --- mDNS-SD traffic reaches every machine on the network by design (Section~\ref{sec:mdns-dnssd}).

Static API keys carry three high-severity threats (\argref{A5}). A stolen key enables spoofing with no time bound because it authenticates the bearer indefinitely. Static keys provide no temporal binding, making session attribution difficult and repudiation easy. Ephemeral credentials expire in ten minutes, limiting the spoofing window to a single rotation cycle. They rotate on a fixed schedule, providing the temporal binding that addresses repudiation. A leaked ephemeral key grants access for at most one cycle. Three medium-severity threats replace the three high-severity ones, all requiring LAN proximity: an attacker on the local network can observe the ephemeral key from mDNS traffic (information disclosure), use that key for up to ten minutes (tampering), or register a rogue \texttt{\_saturn.\_tcp.local.} service to redirect clients (service-level spoofing). Static API keys are exploitable from anywhere on the internet; Saturn's threats require physical or logical presence on the network segment.

The secret leakage data from Meli et al.~\cite{meli2019secrets} makes this comparison concrete. Most leaked API keys are never revoked, and the best automated scanners catch only a fraction of leaks (\argref{A16}). The static key lifecycle fails because keys persist until someone actively removes them, and the evidence shows removal rarely happens. Saturn's ephemeral model sidesteps this failure class: keys expire after ten minutes regardless of whether anyone detects a leak. A Saturn key committed to a public repository would be dead before any scanner could find it. I do not claim that ephemeral credentials make Saturn secure in an enterprise sense. They convert unbounded, internet-scale threats into bounded, LAN-scoped ones --- and the resulting threat surface is documented in the mDNS security literature that Kaiser and Waldvogel~\cite{kaiser2014privacy} established.

\subsection{AP Isolation}

One deployment constraint directly undermines the campus scenario that motivated this work (\argref{A17}). Enterprise and institutional WiFi networks commonly enable AP isolation, which blocks multicast traffic between wireless clients. I confirmed in  that both eduroam and UCSC-Guest block mDNS. Saturn services advertised on these networks are invisible to client devices. This is an infrastructure policy decision and it means Saturn cannot operate on the networks where the access equity motivation is strongest. Saturn works on home networks, office LANs, and lab environments where multicast flows freely. It fails on institutional networks with client isolation.

%----------------------------------------------------------------------
\section{Usability Critique}
\label{sec:usability-critique}

The preceding sections interpret Saturn's functional results and security surface. A single-author inspection against Nielsen's ten usability heuristics~\cite{nielsen1994heuristics}, applied across the three-persona model from Section~\ref{sec:eval-c2}, reveals what those results mean for the people who would use the system.

Saturn's strongest usability property is recognition over recall (\argref{A18}). Discovered services appear as a list with endpoint URLs, backend types, and priority rankings already resolved. Automated assignment of addresses, ports, and failover priorities prevents the configuration slips that manual provisioning invites: malformed URLs, stale endpoints, mismatched ports. This is the mechanism through which Saturn delivers the step reductions the cognitive walkthrough measured. The developer's four remaining steps succeed because none requires recalled knowledge about a provider's API surface.

Three weaknesses emerge, each tied to a design trade-off. First, Saturn's invisibility to end users creates a control gap. Administrators and developers can manipulate service discovery, but end users cannot override misconfigured priorities that degrade their experience. The zero-step end-user experience comes at the cost of zero end-user agency. Second, Saturn's terminology --- ``beacon,'' ``proxy,'' ``service type'' --- matches protocol semantics but diverges from administrator mental models, where ``model endpoint'' is the natural concept. This gap matters less for protocol correctness than for adoption: administrators who cannot map Saturn's vocabulary to their existing infrastructure concepts will resist deploying it. Third, and most critically, Saturn lacks failure diagnostics. When discovery fails, the client reports a generic ``service not available'' error that does not distinguish between AP isolation blocking multicast, no beacon running, or a misconfigured beacon. Nielsen's error recovery heuristic requires that error messages identify the problem and suggest corrective action; Saturn does neither. The AP isolation problem from above compounds this weakness: the most common deployment failure --- institutional WiFi blocking multicast --- produces the least informative error.

\section{Future Work}
\label{sec:future-work}

I chose mDNS-only discovery because it requires no infrastructure beyond the network itself --- no central registry, no DNS records to configure, no server to maintain. That purity is also the protocol's sharpest constraint. The AP isolation barrier (\argref{B5}) demands a hybrid discovery architecture: pairing mDNS with a lightweight unicast fallback that advertises Saturn services via a well-known HTTPS endpoint. The concrete flow would be: a client on eduroam queries \texttt{\_saturn.\_tcp} via mDNS, receives no response, then falls back to \texttt{https://saturn.ucsc.edu/.well-known/saturn} for endpoint metadata. That fallback reintroduces a configuration dependency --- someone must provision and maintain the HTTPS endpoint --- but the dependency falls on the administrator, not the end user. The design goal is to preserve Saturn's zero-configuration property for the student while accepting that institutional networks require institutional infrastructure. Success means a student on eduroam discovers a Saturn service as seamlessly as one on a home network.

Empirical user studies would complement the analytical evaluation. I would run a controlled experiment comparing configuration time and error rates between Saturn and traditional API key provisioning, across participants with varying technical backgrounds, to test whether the step reductions from the cognitive walkthrough translate to real usability gains. Carmona-Galindo et al.~\cite{carmona2025hsi} documented how generative AI adoption varies across student populations at Hispanic-serving institutions; I would prioritize that context to directly test Saturn's potential to reduce the access disparities that Bassignana et al.~\cite{bassignana2025aigap} identified.

Finally, Kaiser and Waldvogel's privacy-preserving mDNS-SD extensions~\cite{kaiser2014efficient} offer a path toward confidential discovery. Their privacy sockets and pairing mechanisms maintain backward compatibility with standard mDNS while encrypting service advertisements to authorized clients. Integrating these extensions would enable Saturn deployments in shared workspaces or medical facilities where broadcasting AI service availability raises privacy concerns the current protocol does not address.

The gap between AI's falling inference costs and its persistent access barriers is a protocol problem, not a cost problem. The infrastructure to close that gap has existed for over a decade; what remained was the evidence that it generalizes and the honesty about where it breaks. If the hybrid discovery and empirical validation work I have outlined succeed, the next generation of students will encounter AI as something the network simply provides.
